\documentclass[a4paper,10pt,notitlepage,oneside]{article}
\usepackage[utf8]{inputenc} % pour les accents
\usepackage[greek,english,french]{babel}



%%%%%%%%%%%%%%%%%%------ packages --------%%%%%%%%%%%%%%%%%%%%%%%%%%%%%%

\usepackage[top=2.5cm, bottom=2.5cm, left=2cm, right=2cm]{geometry}

\usepackage[utf8]{inputenc}
\usepackage{lmodern}
\usepackage{ntheorem}
\usepackage{gensymb}
\usepackage{babel}
\usepackage{dsfont}
\usepackage[dvipsnames]{xcolor}
\usepackage{multicol}
\usepackage{multirow}
\usepackage{fancyhdr}             
\usepackage{float}
\usepackage{amsfonts}
\usepackage{amsmath}
\usepackage{amssymb}
\usepackage{latexsym}
\usepackage{array}
\usepackage{graphicx}
%%%%%%%%%%%%%%-----NEW COMMAND-------%%%%%%%%%%%%%%%%%%%%%%%%%%%%%%%%%%
\newcommand{\dis}{\displaystyle}
\newcommand{\ve}{ \overrightarrow}
\mathchardef\times="2202
\newcommand{\C}{\mathbb{C}}
\newcommand{\R}{\mathbb{R}}
\newcommand{\Q}{\mathbb{Q}}
\newcommand{\Z}{\mathbb{Z}}
\newcommand{\N}{\mathbb{N}}


\newcommand{\HRule}{\rule{\linewidth}{0.5mm}}

\usepackage[cyr]{fourier}
\pagestyle{plain}
%%%%%%%%%%%%%%%%%%%%%%%%%%%%%%%%%%%%%%%%%%%%%%%%%%%
%FONCTION
\newenvironment{norm}{%
\left\Vert
}{%
\right\Vert
}
%%%%%%%%%%%%%%%%%%%%%%%%%%%%%%%%%%%%%%%%%%%%%%%%%%
\AtBeginDocument{
   \renewcommand\tablename{\textsc{Tableau}}
   \renewcommand\figurename{\textsc{Figure}}
}
\usepackage{graphicx}
\usepackage{tabularx} 
\usepackage{multirow}
\usepackage{color}
\definecolor{ENSblue}{RGB}{0,119,139}
\usepackage[backref=true,colorlinks=true,linkcolor=ENSblue,urlcolor=ENSblue]{hyperref}

\usepackage{lipsum}
\usepackage{amsmath}
\usepackage{fancybox}  
\usepackage{listings}
\usepackage[squaren,Gray]{SIunits}
\usepackage{verbatim}
\setcounter{tocdepth}{2}
\usepackage{listings}
\usepackage{float}
\usepackage{microtype}
\usepackage{tcolorbox}
%%%%%%%%%%%%%%%%%%%%%%%%%%%%%%%%%%%%%%%%%%%%%%%%%%%%%%%%%%%%%%%%%%%
%-----------------Header and footer-------------------------

%---EN-TETE-ET-PIED-DE-PAGE----------------------------------------------------
\renewcommand{\headrulewidth}{0.4pt}
\renewcommand{\footrulewidth}{0.4pt}
\setlength{\columnseprule}{0.4pt}% the line between colomn paper

\fancyhead{} % clear all header fields
\fancyhead[LO,CE]{\textbf{\color{BlueViolet}Colles}}
\fancyfoot{} % clear all footer fields
\fancyhead{} % clear all header fields
\fancyhead[LO,CE]{\color{BlueViolet} \textbf{\color{BlueViolet}Colles}}
\fancyfoot{} % clear all footer fields
\fancyfoot[LE,RO]{\color{BlueViolet} \textbf{\color{BlueViolet}Semaine 8}}
\fancyfoot[LO,CE]{\color{BlueViolet} \textbf{\color{BlueViolet}Physique - Chimie}}

\pagestyle{fancy}

%\lhead{fhgfgd}
\cfoot{\thepage}
%%%%%%%%%%%%%%%%%%%%%%%%%%%%%%%%%%%%%%%%%%%%%%%%

\renewcommand{\headrulewidth}{0.4pt}
\renewcommand{\footrulewidth}{0.4pt}
\setlength{\columnseprule}{0.4pt}% the line between colomn paper
% Appliquer la couleur à la ligne de tête et de pied de page
\renewcommand{\headrule}{\color{BlueViolet}\rule{\linewidth}{0.4pt}}  % Ligne de tête en violet
\renewcommand{\footrule}{\color{BlueViolet}\rule{\linewidth}{0.4pt}}  % Ligne de pied de page en violet


%%%%%%%%%%%%%%%%%%%%%%%%%%%%%%%%%%%%%%%%%%%%%%%%
%------- exercices style -----------
\theoremstyle{plain}
\theorembodyfont{\normalfont}
\theoremseparator{~--}
\newtheorem{partie}{\color{BlueViolet} Partie}

\renewcommand{\familydefault}{lmss}

\begin{document}





\begin{center}
    {\color{BlueViolet} \Large \textbf{\color{BlueViolet}Sujet \no 1 - Correction}}
\end{center}

\paragraph{\color{BlueViolet} Exercice 1 : Pouvoir de réchauffement global du méthane}

\begin{enumerate}
    \item Écrire la loi de vitesse. Identifier une constante cinétique apparente.\\
    La réaction est de premier ordre par rapport à chaque réactif ($CH_4$ et $HO^.$), donc la loi de vitesse s'écrit :
    \[
    v = k_0 [CH_4] [HO^.].
    \]
    En supposant que la concentration des radicaux hydroxyles \([HO^.]\) reste constante, on peut définir une constante cinétique apparente \(k\) telle que :
    \[
    k = k_0 [HO^.],
    \]
    ce qui permet de simplifier la loi de vitesse :
    \[
    v = k [CH_4].
    \]
    \begin{equation*}
        \boxed{v = k [CH_4] \quad \text{avec} \quad k = k_0 [HO^.]}
    \end{equation*}

    \item Déterminer la masse restante \(m_{CH_4}(t)\) au bout d’une durée \(t\).\\
    On a, dans un volume donné,
    \begin{equation*}
        v = -\dfrac{d\left[ CH_4 \right]}{dt} = k\left[ CH_4 \right]
    \end{equation*}
    soit
    \begin{equation*}
        \dfrac{d\left[ CH_4 \right]}{dt} + k\left[ CH_4 \right] = 0
    \end{equation*}
    En rappelant la définition de la concentration, $\left[ CH_4 \right](t) = \dfrac{m_{CH_4}(t)}{M_{CH_4}V}$, puis en injectant dans l'équation différentielle, on obtient après simplification :
    \[
    \dfrac{dm_{CH_4}}{dt} = -k m_{CH_4}.
    \]
    En séparant les variables et intégrant entre \(t = 0\) et \(t\), on obtient :
    \[
    \ln \left( \dfrac{m_{CH_4}(t)}{m_0} \right) = -kt.
    \]
    \begin{equation*}
        \boxed{m_{CH_4}(t) = m_0 e^{-kt}}
    \end{equation*}

    \item En déduire la constante cinétique \(k_0\).\\
    Le temps de demi-vie \(t_{1/2}\) est défini par :
    \[
    m_{CH_4}(t_{1/2}) = \dfrac{m_0}{2}.
    \]
    En utilisant la solution obtenue précédemment, on a :
    \[
    \dfrac{1}{2} = e^{-k t_{1/2}} \quad \text{soit} \quad k = \dfrac{\ln 2}{t_{1/2}}.
    \]
    Substituons \(t_{1/2} = 8.2 \, \text{ans}\) :
    \[
    k = \dfrac{\ln 2}{8.2} = 0.0845 \, \text{an}^{-1}.
    \]
    La constante cinétique \(k_0\) est alors calculée à partir de \(k = k_0 [HO^.]\). Sachant que \([HO^.] = 10^{12} \, \text{radicaux/m}^3 = 1.66 \cdot 10^{-12} \, \text{mol/m}^3\), on obtient :
    \begin{equation*}
        \boxed{k_0 = \dfrac{k}{[HO^.]} = \dfrac{0.0845}{1.66 \cdot 10^{-12}} \approx 5.09 \cdot 10^{10} \, \text{m}^3 \cdot \text{mol}^{-1} \cdot \text{an}^{-1}}
    \end{equation*}
    \item On obtient :
    \begin{equation*}
        I_{CH_4}(\tau) = \displaystyle{\int_0^\tau}e^{-kt}dt = \left[-\dfrac{1}{k}e^{-kt}\right]_0^\tau \quad \text{soit} \quad \boxed{I_{CH_4}(\tau) = \dfrac{1}{k}\left(1-e^{-kt}\right)}
    \end{equation*}
    \item Le méthane se décompose plus vite que le $CO_2$, son effet est donc plus important sur des durées courtes que sur des durées longues, d’où l’allure décroissante. Dire qu’une tonne de méthane équivaut à $25$ tonnes de $CO_2$ est abusif : ce n’est vrai qu’à un horizon d’environ un siècle. Si on se place à un horizon plus proche, par exemple 20 ans, une tonne de méthane a le même effet sur le réchauffement à cette échéance que 80 tonnes de $CO_2$.  
\end{enumerate}


\\

\paragraph{\color{BlueViolet} Exercice 2 : Décomposition de l'éthanal}

\begin{enumerate}
    \item Nommer les espèces et écrire l'équation de réaction.\\
    
    Les espèces présentes sont :
    \begin{itemize}
        \item \textbf{\color{BlueViolet}Réactif :} l'éthanal $\text{CH}_3\text{CHO}$,
        \item \textbf{\color{BlueViolet}Produits :} le méthane $\text{CH}_4$ et le monoxyde de carbone $\text{CO}$.
    \end{itemize}
    L'équation de la réaction s'écrit :
    \begin{equation*}
        \boxed{\text{CH}_3\text{CHO}(g) \longrightarrow \text{CH}_4(g) + \text{CO}(g)}
    \end{equation*}
    
    \item Construire le tableau d'avancement à l'instant $t$ en fonction de l'avancement $\xi(t)$.\\
    
    Le tableau d'avancement s'écrit :
    \[
    \begin{array}{c|c|c|c}
        \text{État} & \text{CH}_3\text{CHO} & \text{CH}_4 & \text{CO} \\
        \hline
        t = 0 & n_0 & 0 & 0 \\
        t & n_0 - \xi & \xi & \xi \\
    \end{array}
    \]
    
    \item Montrer que l'on peut suivre l'avancement par la mesure d'une seule grandeur physique.\\
    La quantité totale de gaz dans le milieu réactionnel est donnée par :
    \[
        n_{\text{gaz}} = n_{\text{CH}_3\text{CHO}} + n_{\text{CH}_4} + n_{\text{CO}} = n_0 + \xi
    \]
    La quantité totale varie donc avec l'avancement de la réaction. En appliquant la loi des gaz parfaits, la pression totale vaut :
    \[
        p = n_{\text{gaz}} \dfrac{RT}{V} = \left( n_0 + \xi \right) \dfrac{RT}{V} = p_0 + \dfrac{RT}{V} \xi
    \]
    Ainsi, la mesure de la pression $p$ permet de déterminer l'avancement $\xi$.

    \item On constate expérimentalement que la fonction 
    \[
    F(t) = \dfrac{- (p(t) - p_0)}{p(t) - 2p_0}
    \]
    est proportionnelle à $t$. Montrer qu'une réaction d'ordre 2 est compatible avec ces résultats.\\
    
    La vitesse d'une réaction d'ordre 2 est donnée par :
    \[
        v = -\dfrac{d[\text{CH}_3\text{CHO}]}{dt} = k[\text{CH}_3\text{CHO}]^2
    \]
    En séparant les variables et en intégrant :
    \[
        \displaystyle{\int_{n_0/V}^{(n_0 - \xi)/V}} \dfrac{d[\text{CH}_3\text{CHO}]}{[\text{CH}_3\text{CHO}]^2} = -k \displaystyle{\int_{0}^{t}} dt
    \]
    Après intégration :
    \[
        -\dfrac{1}{n_0 - \xi} + \dfrac{1}{n_0} = -kt
    \]
    Avec $n_0 = \dfrac{p_0 V}{RT}$ et $\xi = \dfrac{V}{RT}(p - p_0)$, on obtient :
    \[
        \boxed{F(t) = \dfrac{- (p - p_0)}{p - 2p_0} = -\dfrac{n_0 k t}{V}}
    \]
    Cette relation montre que la fonction $F(t)$ est bien proportionnelle au temps $t$, confirmant que l'ordre 2 est compatible.

    \item Calculer le temps de demi-réaction.\\
    
    Au temps de demi-réaction $t_{1/2}$, il reste la moitié de la quantité initiale de $\text{CH}_3\text{CHO}$. Ainsi :
    \[
        n_{\text{gaz}}(t_{1/2}) = n_0 + \dfrac{n_0}{2} = \dfrac{3}{2} n_0 \implies p(t_{1/2}) = \dfrac{3}{2} p_0
    \]
    En substituant dans la relation précédemment établie, le temps de demi-réaction est donné par :
    \[
        \boxed{t_{1/2} = \dfrac{V}{n_0 k}}
    \]
\end{enumerate}


\newpage
\begin{center}
    {\color{BlueViolet} \Large \textbf{\color{BlueViolet}Sujet 2 - Correction}}
\end{center}


\paragraph{\color{BlueViolet} Exercice 1 : Estimation de l'âge de la Terre par la désintégration de l'uranium}

\begin{enumerate}
    \item Écrire la loi d'évolution du nombre d’atomes de \(^{238}U\) en fonction du temps.\\
    La loi de désintégration radioactive est donnée par :
    \[
    \dfrac{dN}{dt} = -\lambda N.
    \]
    En séparant les variables et en intégrant entre \(t = 0\) (où \(N = N_0\)) et un temps \(t\), on obtient :
    \begin{equation*}
        \boxed{N(t) = N_0 e^{-\lambda t}}
    \end{equation*}

    \item Exprimer le nombre d’atomes désintégrés de \(^{238}U\) au cours du temps.\\
    Le nombre d'atomes désintégrés au temps \(t\) est donné par :
    \[
    N_{\text{désintégrés}} = N_0 - N(t).
    \]
    En utilisant la relation \(N(t) = N_0 e^{-\lambda t}\), :
    \begin{equation*}
        \boxed{N_{\text{désintégrés}} =N_0 - N_0 e^{-\lambda t} = N_0 (1 - e^{-\lambda t})}
    \end{equation*}

    \item Relation entre \(N_{U}\), \(N_{Pb}\) et \(t\).\\
    Tous les atomes de \(^{206}Pb\) proviennent de la désintégration de \(^{238}U\), donc :
    \[
    N_{Pb} = N_0 - N_{U},
    \]
    où \(N_0\) est le nombre initial d'atomes de \(^{238}U\). En utilisant \(N_{U} = N_0 e^{-\lambda t}\), on a :
    \[
    N_{Pb} = N_0 (1 - e^{-\lambda t}).
    \]
    En remplaçant \(N_0\) par \(N_{U} / e^{-\lambda t}\), on obtient :
    \[
    e^{\lambda t} = 1 + \dfrac{N_{Pb}}{N_{U}}.
    \]
    En prenant le logarithme :
    \[
    t = \dfrac{1}{\lambda} \ln \left( 1 + \dfrac{N_{Pb}}{N_{U}} \right).
    \]
    \begin{equation*}
        \boxed{t = \dfrac{1}{\lambda} \ln \left( 1 + \dfrac{N_{Pb}}{N_{U}} \right)}
    \end{equation*}

    \item Calcul de l’âge de l’échantillon.\\
    La constante radioactive \(\lambda\) est liée à la période de demi-vie \(t_{1/2}\) par :
    \[
    \lambda = \dfrac{\ln 2}{t_{1/2}}.
    \]
    En remplaçant \(t_{1/2} = 4.5 \cdot 10^9 \, \text{ans}\), on obtient :
    \[
    \lambda = \dfrac{\ln 2}{4.5 \cdot 10^9} = 1.54 \cdot 10^{-10} \, \text{an}^{-1}.
    \]
    Substituons dans l’expression de \(t\) :
    \[
    t = \dfrac{1}{1.54 \cdot 10^{-10}} \ln \left( 1 + \dfrac{6 \cdot 10^{12}}{2 \cdot 10^{12}} \right).
    \]
    Le rapport \(\dfrac{N_{Pb}}{N_{U}} = \dfrac{6}{2} = 3\), donc :
    \[
    t = \dfrac{1}{1.54 \cdot 10^{-10}} \ln(4).
    \]
    Puisque \(\ln(4) = 2 \ln(2)\), on a :
    \[
    t = \dfrac{1}{1.54 \cdot 10^{-10}} \cdot 2 \cdot 0.693 = 9.0 \cdot 10^9 \, \text{ans}.
    \]
    \begin{equation*}
        \boxed{t = 9.0 \cdot 10^9 \, \text{ans}}
    \end{equation*}

    Ainsi, l’âge de l’échantillon est d’environ 9 milliards d’années.
\end{enumerate}



\paragraph{\color{BlueViolet} Exercice 2 : Modélisation d'un bâtiment sous séisme}

\begin{enumerate}
    \item En appliquant le PFD le long de l'axe de déplacement, montrer que l'équation différentielle du système peut se mettre sous la forme
    \begin{equation*}
        \dfrac{d^2x_1}{dt^2}+2\xi\omega_0\dfrac{dx_1}{dt}+\omega_0^2x_1(t) = 0
    \end{equation*}
    et identifier $\omega_0$ la pulsation propre et $\xi$ le taux d'amortissement.\\
    
    On peut directement écrire 
    \begin{equation*}
        m_1\dfrac{d^2x_1}{dt^2}+\eta_1\dfrac{dx_1}{dt}+k_1x_1(t) = 0
    \end{equation*}
    Par identification,
    \begin{equation*}
        \boxed{\xi = \dfrac{\eta_1}{2\sqrt{k_1m_1}}}\quad 
        \boxed{\omega_0 = \sqrt{\dfrac{k_1}{m_1}}}
    \end{equation*}
    
    \item Quelle est l'expression littérale du coefficient d'amortissement $\eta_{\text{cr}}$ au-delà duquel il n'y aura pas d'oscillation ? (régime apériodique, ou régime sur-critique).\\

    Ce coefficient est celui qui vérifie $\xi>1$, donc $\boxed{\eta_{cr}=2\sqrt{k_1m_1}}$
    
    \item En règle générale, les bâtiments réagissent en régime pseudo-périodique. Quelle est la forme de la solution dans ce cas-là? On prendra $x_1(t=0)=0$ et $\dfrac{dx_1}{dt}(t=0) = 0$. On notera également $\omega_D = \omega_0\sqrt{1-\xi^2}$ la pulsation propre du système amorti.\\
    
    Dans ce cas on aura
    \begin{equation*}
        \boxed{x_1(t) = e^{-\xi\omega_0}\left(A\cos\left(\omega_D.t\right)+B\sin\left(\omega_D.t\right)\right)}
    \end{equation*}
    
    \item On s'intéresse maintenant à la modélisation de deux étages. Le deuxième étage voit sa position sur l'axe transversal $x_2(t)$, et il est lié à l'étage précédent par une structure qui se comporte comme un ressort de rigidité $k_2$, de longueur à vide nulle, et d'amortissement $\eta_2$.
    Quelle est l'équation différentielle qui régit la position du premier étage?\\
    La seule différence avec la partie précédente concerne l'effort supplémentaire qui s'exerce sur la masse $m_1$ liée au déplacement différentiel entre $m_1$ et $m_2$.
    \begin{equation*}
        \boxed{
         m_1\dfrac{d^2x_1}{dt^2}+\eta_1\dfrac{dx_1}{dt}+k_1 x_1(t) +\textcolor{red}{k_2\left(x_1-x_2\right)} = 0
        }
    \end{equation*}
    \item Quelle est l'équation différentielle qui régit la position du deuxième étage?\\
    Les seules forces qui s'exercent sur la masse $m_2$ sont le poids, l'amortissement, et le rappel du ressort qui lie $m_1$ à $m_2$.
    \begin{equation*}
        \boxed{
         m_2\dfrac{d^2x_2}{dt^2}+\eta_2\dfrac{dx_2}{dt}+ k_2\left(x_2-x_1\right) = 0
        }
    \end{equation*}
    \item Exprimer l'équation matricielle à résoudre.\\
    On aura ici le système
    \begin{equation*}
        \boxed{
        M .\ddot X + C .\dot X+K.\dot X = 0
        }
    \end{equation*}
    où $0$ est ici la matrice nulle.
\end{enumerate}


\newpage


\begin{center}
    {\color{BlueViolet} \Large \textbf{\color{BlueViolet}Sujet \no 3 - Correction}}
\end{center}

\paragraph{\color{BlueViolet} Exercice 1 : Transport du méthane}

\begin{enumerate}
    \item Exprimer la quantité de matière totale de gaz $n_{\text{gaz},f}$ dans le mélange en fonction de $n$ et $\xi_f$.\\
    D'après le bilan de matière dans l'état final :
    \[
    n_{\text{gaz},f} = n_{CH_4,f} + n_{\ce{H2},f} = (n - \xi_f) + 2 \xi_f = n + \xi_f.
    \]
    \begin{equation*}
        \boxed{n_{\text{gaz},f} = n + \xi_f}
    \end{equation*}

    \item Montrer que les pressions partielles en méthane et en dihydrogène sont données par : 
    \[
    p_{CH_4} = \dfrac{1 - \alpha}{1 + \alpha} P \quad \text{et} \quad p_{\ce{H2}} = \dfrac{2\alpha}{1 + \alpha} P
    \]
    avec $\alpha = \dfrac{\xi_f}{n}$.\\
    Les pressions partielles sont données par :
    \[
    p_{CH_4} = \dfrac{n_{CH_4,f}}{n_{\text{gaz},f}} P = \dfrac{n - \xi_f}{n + \xi_f} P = \dfrac{1 - \alpha}{1 + \alpha} P,
    \]
    \[
    p_{\ce{H2}} = \dfrac{n_{\ce{H2},f}}{n_{\text{gaz},f}} P = \dfrac{2 \xi_f}{n + \xi_f} P = \dfrac{2\alpha}{1 + \alpha} P.
    \]
    \begin{equation*}
        \boxed{p_{CH_4} = \dfrac{1 - \alpha}{1 + \alpha} P \quad \text{et} \quad p_{\ce{H2}} = \dfrac{2\alpha}{1 + \alpha} P}
    \end{equation*}

    \item Déterminer la valeur de $\alpha$. Commenter le résultat.\\
    La constante d'équilibre est donnée par :
    \[
    K^\circ = \dfrac{(p_{\ce{H2}})^2}{p_{CH_4} p^\circ},
    \]
    ce qui donne en remplaçant $p_{CH_4}$ et $p_{\ce{H2}}$ :
    \[
    K^\circ = \dfrac{\left(\dfrac{2\alpha}{1+\alpha} P\right)^2}{\dfrac{1-\alpha}{1+\alpha} P \cdot p^\circ} = \dfrac{4\alpha^2 P}{p^\circ(1-\alpha)}.
    \]
    Pour $K^\circ \ll 1$, on suppose $\alpha \ll 1$ et on approxime $(1-\alpha) \approx 1$. On obtient alors :
    \[
    \alpha = \sqrt{\dfrac{K^\circ p^\circ}{4P}}.
    \]
    En remplaçant les valeurs numériques :
    \[
    \alpha = \sqrt{\dfrac{4 \cdot 10^{-9} \cdot 1 \, \text{bar}}{4 \cdot 70 \, \text{bar}}} = 2 \cdot 10^{-6}.
    \]
    \begin{equation*}
        \boxed{\alpha = 2 \cdot 10^{-6}}
    \end{equation*}
    Ce résultat indique que la dissociation du méthane est très faible, ce qui est souhaitable pour limiter les pertes de méthane dans le gazoduc.

    \item Déterminer la masse de carbone solide qui se dépose chaque fois qu’une tonne de méthane transite dans le gazoduc.\\
    D'après le bilan de matière, $n_{\ce{C}} = \xi_f = \alpha n$. La quantité de matière initiale de méthane est donnée par :
    \[
    n = \dfrac{m_{CH_4}}{M_{CH_4}} = \dfrac{10^6 \, \text{g}}{16 \, \text{g} \cdot \text{mol}^{-1}} = 6.25 \cdot 10^4 \, \text{mol}.
    \]
    La masse de carbone solide déposée est alors :
    \[
    m_{\ce{C}} = n_{\ce{C}} M_{\ce{C}} = \alpha n \cdot M_{\ce{C}} = 2 \cdot 10^{-6} \cdot 6.25 \cdot 10^4 \cdot 12 = 1.5 \, \text{g}.
    \]
    \begin{equation*}
        \boxed{m_{\ce{C}} = 1.5 \, \text{g}}
    \end{equation*}
\end{enumerate}



\paragraph{\color{BlueViolet} Exercice 2 : Un glaçon dans un pastis}
\begin{enumerate}
    \item Faire un schéma légendé de la situation\\
        
    \begin{figure}[H]
        \centering
        \includegraphics[width=0.4\linewidth]{S9/glaçon2.png}
    \end{figure}

    \item Calculer, pour une position $z$, le volume immergé et le volume émergé du glaçon.\\
    Le volume immergé vaut
    \begin{equation*}
        \boxed{
        V_{\text{immergé}} = \left(\dfrac{l}{2}-z(t)\right)Ll
        }
    \end{equation*}
    Et le volume émergé vaut
    \begin{equation*}
        \boxed{
        V_{\text{émergé}} = \left(\dfrac{l}{2}+z(t)\right)Ll
        }
    \end{equation*}
    
    \item L'équilibre est atteint lorsque le poids équilibre la poussée d'Archimède. Donner l'expression littérale puis la valeur numérique de la position d'équilibre $z_{eq}$.\\
    Le poids s'exprime, quelle que soit la position, 
    \begin{equation*}
        \overrightarrow{P} = -\rho_{\text{glace}}V.\overrightarrow{e_z} =  -\rho_{\text{glace}}l^2Lg.\overrightarrow{e_z}
    \end{equation*}
    La poussée d'Archimède dépend de la position du glaçon, elle s'exprime :
    
    \begin{equation*}
        \overrightarrow{\Pi} = \rho_{\text{anis}}V_{\text{immergé}}g.\overrightarrow{e_z} = \rho_{\text{anis}} \left(\dfrac{l}{2}-z(t)\right)Llg.\overrightarrow{e_z}
    \end{equation*}
    Lors de l'équilibre,
    \begin{equation*}
        \overrightarrow{P} + \overrightarrow{\Pi} = \overrightarrow{0}
    \end{equation*}
    donc
    \begin{equation*}
    \begin{split}
        &\rho_{\text{glace}}l^2Lg = \rho_{\text{anis}} \left(\dfrac{l}{2}-z_{eq}\right)Llg\\
        \implies& \boxed{z_{eq} = \left(\dfrac{1}{2}-\dfrac{\rho_{\text{glace}}}{\rho_{\text{anis}}}\right)l}
    \end{split}
    \end{equation*}
    \item En négligeant les frottements de l'air et la poussée d'Archimède dans l'air, établir l'équation différentielle vérifiée par la position du glaçon.\\

    Le PFD donne 
    \begin{equation*}
        \rho_{\text{glace}}l^2L\dfrac{d^2z}{dt^2} = - \rho_{\text{glace}}l^2Lg+\rho_{\text{anis}}Llg\left(\dfrac{l}{2}-z\right)-\eta Ll\dfrac{dz}{dt}
    \end{equation*}
    Ce qui, après manipulation, donne
    \begin{equation*}
    \boxed{
       \dfrac{d^2z}{dt^2} +\dfrac{\eta}{ \rho_{\text{glace}}l}\dfrac{dz}{dt}+\dfrac{\rho_{\text{anis}}}{\rho_{\text{glace}}}\dfrac{g}{l}z= g\left(\dfrac{1}{2}\dfrac{\rho_{\text{anis}}}{\rho_{\text{glace}}}-1\right)
       }
    \end{equation*}
    
    
    \item Vérifier que l'on retrouve bien l'expression de la position d'équilibre.\\

    A l'équilibre
    \begin{equation*}
       \dfrac{\rho_{\text{anis}}}{\rho_{\text{glace}}}\dfrac{g}{l}z_{eq}= g\left(\dfrac{1}{2}\dfrac{\rho_{\text{anis}}}{\rho_{\text{glace}}}-1\right)
    \end{equation*}
    soit 
    \begin{equation*}
        \boxed{
        z_{eq} = \left(\dfrac{1}{2}-\dfrac{\rho_{\text{glace}}}{\rho_{\text{anis}}}\right)l
        }
    \end{equation*}
    \item Réécrire l'équation différentielle en exprimant le second membre en fonction de $z_{eq}$.\\

    
    \begin{equation*}
    \boxed{
       \dfrac{d^2z}{dt^2} +\dfrac{\eta}{ \rho_{\text{glace}}l}\dfrac{dz}{dt}+\dfrac{\rho_{\text{anis}}}{\rho_{\text{glace}}}\dfrac{g}{l}z= \dfrac{\rho_{\text{anis}}}{\rho_{\text{glace}}}\dfrac{g}{l}z_{eq}
       }
    \end{equation*}
    
    \item En posant $Z = z-z_{eq}$ Mettre cette équation sous la forme canonique.\\

    On peut remarquer que $\dfrac{dZ}{dt} = \dfrac{dz}{dt}-\dfrac{dz_{eq}}{dt} = \dfrac{dz}{dt}$ (de même pour $\dfrac{d^2Z}{dt^2}$). Donc l'équation différentielle devient :
    
    \begin{equation*}
    \begin{split}
       & \dfrac{d^2Z}{dt^2} +\dfrac{\eta}{ \rho_{\text{glace}}l}\dfrac{dZ}{dt}+\dfrac{\rho_{\text{anis}}}{\rho_{\text{glace}}}\dfrac{g}{l}\left(z-z_{eq}\right) = 0\\
        \implies&\boxed{ \dfrac{d^2Z}{dt^2} +\dfrac{\eta}{ \rho_{\text{glace}}l}\dfrac{dZ}{dt}+\dfrac{\rho_{\text{anis}}}{\rho_{\text{glace}}}\dfrac{g}{l}Z = 0 }
    \end{split}
    \end{equation*}
    Par identification,
    \begin{equation*}
        \boxed{Q = \dfrac{1}{\eta}\sqrt{ \rho_{\text{glace}} \rho_{\text{anis}}lg}} \quad \text{et} \quad \boxed{\omega_0 = \sqrt{\dfrac{ \rho_{\text{anis}}g}{ \rho_{\text{glace}}l}}}
    \end{equation*}
    \item Le régime est-il apériodique, critique ou pseudo-périodique ? On utilisera les valeurs numériques données dans l'énoncé.
    Avec les valeurs numériques, on obtient $Q = $. On est donc en régime pseudo périodique.
     \item Quelle est l'équation du mouvement du glaçon ?
     L'équation du glaçon est de la forme 
     \begin{equation*}
     \boxed{
         Z(t) = e^{\frac{-t}{\tau}}\left(A\cos\left(w_P.t\right)+B\sin\left(w_P.t\right)\right) \quad A,B \in \R
         }
     \end{equation*}
     avec 
     \begin{equation*}
         w_P = \sqrt{1-\dfrac{1}{4Q^2}}
     \end{equation*}


\end{document}





